\chapter{A Matemática dos Sistemas Biológicos}
\label{cap:04}

A matemática é a linguagem com que a Biologia de Sistemas formula seus argumentos mais
precisos. Enquanto a descrição verbal de um processo biológico pode ser ambígua, uma
equação diferencial é inequívoca: especifica exatamente como as variáveis de estado
evoluem no tempo, dado um conjunto de parâmetros e condições iniciais. Mas a matematização
da biologia não é um exercício de formalismo vazio --- é uma ferramenta para prever,
distinguir hipóteses e projetar experimentos.

Este capítulo desenvolve as ferramentas matemáticas indispensáveis para a modelagem de
sistemas dinâmicos biológicos. Começamos pelos modelos em tempo discreto --- o mapa
logístico e as cadeias de Markov --- e avançamos para os sistemas de equações diferenciais
ordinárias (EDOs), cobrindo retratos de fase, análise de estabilidade, bifurcações, ciclos
limite e análise de sensibilidade. O fio condutor é sempre a pergunta biológica: como a
topologia matemática de um sistema se relaciona com seu comportamento funcional?

\section{Modelos Discretos e Recursivos}

\begin{definicaoenv}[Sistema Dinâmico em Tempo Discreto]
Um sistema dinâmico em tempo discreto é aquele em que o tempo avança em passos inteiros
$t = 0, 1, 2, \ldots$ A dinâmica é descrita por uma relação de recorrência
\begin{equation}
  x_{t+1} = f(x_t),
\end{equation}
onde $f$ é uma função que mapeia o estado atual $x_t$ no estado do próximo instante.
\end{definicaoenv}

Muitos processos biológicos ocorrem naturalmente em etapas discretas: gerações
não-sobrepostas de insetos sazonais, divisão celular sincronizada em culturas, ciclos de
replicação viral, ou simplesmente coletas experimentais em pontos fixos de tempo. Para esses
sistemas, a modelagem discreta é mais natural do que as equações diferenciais contínuas.

\subsection{O Mapa Logístico e o Caos Determinístico}

Em 1976, o ecologista Robert May publicou na \textit{Nature} um artigo que se tornaria
clássico: ``\textit{Simple mathematical models with very complicated dynamics}''. May
demonstrava que uma equação de uma única variável --- o mapa logístico --- era capaz de
exibir um leque de comportamentos que ia da extinção ao equilíbrio estável, passando por
oscilações periódicas e chegando ao caos determinístico. O mapa logístico é definido por
\begin{equation}
  x_{t+1} = r\, x_t(1 - x_t),
\end{equation}
onde $x_t \in [0,1]$ representa a densidade populacional normalizada e $r > 0$ é a taxa de
crescimento intrínseca.

O comportamento do sistema depende criticamente de $r$. Para $0 < r < 1$, a população se
extingue: $x_t \to 0$. Para $1 < r < 3$, o sistema converge para o ponto fixo estável
$x^* = 1 - 1/r$. Quando $r$ ultrapassa 3, emerge um ciclo estável de período~2. Aumentando
$r$ além de aproximadamente 3{,}44, o período dobra novamente para~4, depois para~8, num
processo de \emph{duplicação de período} (Feigenbaum, 1978) que conduz ao caos para
$r > 3{,}57$. Na região caótica, trajetórias iniciadas em condições iniciais arbitrariamente
próximas divergem exponencialmente --- sensibilidade às condições iniciais que Lorenz chamou
de ``efeito borboleta''.

A estabilidade de um ponto fixo $x^*$ (onde $f(x^*) = x^*$) é determinada pela derivada de
$f$ nesse ponto. Expandindo em série de Taylor em torno de $x^*$, a perturbação
$\xi_t = x_t - x^*$ evolui segundo
\begin{equation}
  \xi_{t+1} \approx f'(x^*)\,\xi_t
  \quad\Rightarrow\quad
  \xi_t \approx [f'(x^*)]^t\,\xi_0.
\end{equation}
O ponto fixo é \textbf{estável} se $|f'(x^*)| < 1$ e \textbf{instável} se $|f'(x^*)| > 1$.
Para o ponto não-trivial $x^* = 1 - 1/r$, tem-se $f'(x^*) = 2 - r$, de modo que a
estabilidade é garantida para $|2 - r| < 1$, ou seja, $1 < r < 3$.

\begin{figure}[htbp]
\centering
\begin{tikzpicture}[scale=0.95]
    % Curva de duplicação de período (estilizada)
    \draw[->, thick] (0,0) -- (10,0) node[right, font=\small] {$r$};
    \draw[->, thick] (0,0) -- (0,5) node[above, font=\small] {$x^*$};

    % Regime 0 < r < 1 (extinção)
    \draw[thick, azulescuro] (0.5,0.5) -- (2.0,0.5);
    \node[above, font=\footnotesize] at (1.25,0.5) {Extinção};

    % Regime 1 < r < 3 (ponto fixo)
    \draw[thick, verdeacido] (2.0,2.0) -- (5.5,1.0);
    \node[above, font=\footnotesize] at (3.7,1.7) {Ponto fixo estável};

    % Bifurcação de período 2
    \draw[thick, laranjacel] (5.5,1.0) -- (5.5,2.5);
    \draw[thick, laranjacel] (5.5,1.0) .. controls (6.5,0.5) and (7.5,0.3) .. (8.0,0.7);
    \draw[thick, laranjacel] (5.5,2.5) .. controls (6.5,2.8) and (7.5,3.0) .. (8.0,2.5);
    \node[right, font=\footnotesize, laranjacel] at (8.0,1.6) {Período 2};

    % Bifurcação 4, 8, caos
    \draw[thick, roxodna] (8.0,0.7) -- (8.2,0.5);
    \draw[thick, roxodna] (8.0,2.5) -- (8.2,2.7);
    \draw[thick, roxodna] (8.0,0.3) -- (8.4,0.8);
    \draw[thick, roxodna] (8.0,2.7) -- (8.4,2.4);
    \draw[thick, roxodna, dashed] (8.4,1.5) -- (9.5,1.5);
    \node[right, font=\footnotesize, roxodna] at (9.5,1.5) {Caos};

    % Marcas
    \draw[dashed] (5.5,0) -- (5.5,2.5);
    \node[below, font=\footnotesize] at (5.5,-0.05) {$r=3$};
    \draw[dashed] (8.4,0) -- (8.4,1.5);
    \node[below, font=\footnotesize] at (8.4,-0.05) {$r \approx 3{,}57$};
\end{tikzpicture}
\caption{Diagrama de bifurcações do mapa logístico $x_{t+1} = r x_t(1 - x_t)$. Para $0 < r < 1$, a população se extingue ($x^* = 0$, linha azul). Para $1 < r < 3$, o sistema converge a um ponto fixo estável $x^* = 1 - 1/r$ (curva verde). Em $r = 3$, ocorre a primeira \emph{bifurcação de duplicação de período}: o ponto fixo se torna instável e surge um ciclo de período 2 (laranja). Sucessivas duplicações de período (4, 8, 16, ...) culminam em caos determinístico para $r \gtrsim 3{,}57$ (região roxa).}
\label{fig:04-logistico}
\end{figure}

\begin{lstlisting}[language=Python, caption=Diagrama de bifurcacao do mapa logistico]
import numpy as np
import matplotlib.pyplot as plt

r_vals = np.linspace(2.5, 4.0, 1000)
x = 0.5 * np.ones(len(r_vals))

for _ in range(500):          # transiente
    x = r_vals * x * (1 - x)

r_plot, x_plot = [], []
for _ in range(200):          # regime permanente
    x = r_vals * x * (1 - x)
    r_plot.extend(r_vals)
    x_plot.extend(x)

plt.figure(figsize=(10, 5))
plt.plot(r_plot, x_plot, ',k', alpha=0.2)
plt.xlabel('r'); plt.ylabel('x*')
plt.title('Diagrama de bifurcacao --- mapa logistico')
plt.tight_layout(); plt.show()
\end{lstlisting}

\subsection{Processos Estocásticos e Cadeias de Markov}

Nem todo processo biológico é determinístico. Quando o número de moléculas ou indivíduos é
pequeno, as flutuações estocásticas tornam-se relevantes e podem alterar qualitativamente a
dinâmica. Um framework natural para modelar esse regime é o das \emph{cadeias de Markov},
nas quais o estado futuro do sistema depende apenas do estado presente --- a chamada
\emph{propriedade de Markov} --- e não de toda a história passada.

\begin{definicaoenv}[Cadeia de Markov em Tempo Discreto]
Uma cadeia de Markov em tempo discreto é definida por um conjunto de estados
$\{s_1, s_2, \ldots, s_n\}$ e por uma \emph{matriz de transição} $\mathbf{P}$, onde
$P_{ij}$ é a probabilidade de transitar do estado $s_i$ para o estado $s_j$ num único passo.
A distribuição de probabilidade dos estados evolui como
\begin{equation}
  \boldsymbol{\pi}_{t+1} = \mathbf{P}^T \boldsymbol{\pi}_t,
\end{equation}
onde $\boldsymbol{\pi}_t$ é o vetor de probabilidades no instante $t$.
\end{definicaoenv}

No limite de longo prazo, uma cadeia de Markov \emph{ergódica} (irredutível e aperiódica)
converge para a \emph{distribuição estacionária} $\boldsymbol{\pi}^*$, que satisfaz
$\boldsymbol{\pi}^* = \mathbf{P}^T \boldsymbol{\pi}^*$ --- ou seja, $\boldsymbol{\pi}^*$
é o autovetor de $\mathbf{P}^T$ associado ao autovalor 1.

Um exemplo biológico paradigmático é o modelo de \emph{canal iônico} de dois estados:
Fechado ($F$) e Aberto ($A$). A matriz de transição é
\begin{equation}
  \mathbf{P} = \begin{bmatrix} 1 - \alpha & \alpha \\ \beta & 1 - \beta \end{bmatrix},
\end{equation}
onde $\alpha$ é a probabilidade de abertura por passo e $\beta$ a de fechamento. A
distribuição estacionária é obtida resolvendo $\pi_F \alpha = \pi_A \beta$:
\begin{equation}
  \pi_A = \frac{\alpha}{\alpha + \beta}, \quad \pi_F = \frac{\beta}{\alpha + \beta}.
\end{equation}
Assim, a fração de canais abertos no equilíbrio é determinada apenas pela razão das taxas
de abertura e fechamento --- resultado que conecta o nível molecular (transições individuais)
ao nível macroscópico (condutância elétrica média da membrana).

Em biologia, cadeias de Markov descrevem também transições entre estados conformacionais de
proteínas, dinâmica de expressão gênica em células individuais, e modelos epidemiológicos
estocásticos (SIR em populações finitas), nos quais a extinção estocástica do vírus pode
ocorrer mesmo quando o modelo determinístico prevê persistência.


\section{Sistemas de EDOs e Retratos de Fase}

A transição dos modelos discretos para os contínuos é natural quando o número de
indivíduos é grande e o tempo pode ser tratado como variável contínua. O
\emph{sistema de equações diferenciais ordinárias} (EDOs) é a linguagem padrão para
modelar a dinâmica contínua de sistemas biológicos:
\begin{equation}
  \frac{d\mathbf{x}}{dt} = \mathbf{f}(\mathbf{x},\, \boldsymbol{\theta}),
\end{equation}
onde $\mathbf{x}(t) \in \mathbb{R}^n$ é o vetor de variáveis de estado (concentrações,
populações, potenciais elétricos) e $\boldsymbol{\theta}$ é o vetor de parâmetros.
O lado direito $\mathbf{f}(\mathbf{x})$ define um \emph{campo vetorial} que associa a
cada ponto do espaço de estados um vetor velocidade --- a direção e magnitude em que o
sistema evolui instantaneamente.

O \emph{retrato de fase} é a representação geométrica das trajetórias no espaço de
estados, independentemente do tempo. Em sistemas bidimensionais, ele revela de forma
imediata a natureza qualitativa da dinâmica --- convergência para equilíbrio, oscilações
periódicas, trajetórias heteroclínicas --- sem exigir solução analítica das EDOs.

As \emph{nulclinas} são curvas no espaço de fase onde uma derivada específica se anula:
a \emph{$x$-nulclina} é o conjunto $\{\dot{x} = 0\}$ (onde o movimento é puramente
vertical) e a \emph{$y$-nulclina} é o conjunto $\{\dot{y} = 0\}$ (movimento puramente
horizontal). Os \emph{pontos de equilíbrio} --- onde $\mathbf{f}(\mathbf{x}^*) = \mathbf{0}$
--- estão sempre nas interseções das nulclinas. Construir o retrato de fase exige: (1)
calcular e plotar as nulclinas; (2) identificar equilíbrios; (3) desenhar o campo de
direções em cada região; (4) integrar trajetórias numericamente.

\subsection{O Sistema de Lotka-Volterra}

O sistema predador-presa de Lotka-Volterra, proposto independentemente por Alfred Lotka
(1920) e Vito Volterra (1926), é um dos modelos de EDOs mais estudados em biologia:
\begin{align}
  \frac{dx}{dt} &= \alpha x - \beta xy \\[4pt]
  \frac{dy}{dt} &= \delta xy - \gamma y,
\end{align}
onde $x$ é a população de presas (crescimento intrínseco $\alpha$, predação $\beta$) e
$y$ é a população de predadores (eficiência de conversão $\delta$, mortalidade $\gamma$).
O sistema possui dois pontos de equilíbrio: a origem $(0,0)$ --- extinção de ambas as
espécies --- e o ponto não-trivial $(x^*, y^*) = (\gamma/\delta,\, \alpha/\beta)$.

A análise de estabilidade desse ponto de equilíbrio, via linearização (ver Seção~4.3),
revela que os autovalores do jacobiano são puramente imaginários, indicando oscilações
conservativas, não amortecidas. As trajetórias formam órbitas fechadas no plano de fase:
as populações oscilam indefinidamente, sem convergir para um equilíbrio fixo nem divergir.
Dados históricos de lince e lebre do Canadá (Elton e Nicholson, 1942) mostram oscilações
qualitativamente consistentes com esse modelo.

As nulclinas do sistema de Lotka-Volterra têm interpretação biológica clara: a
$x$-nulclina $y = \alpha/\beta$ (linha horizontal) indica o número de predadores que
mantém as presas em equilíbrio; a $y$-nulclina $x = \gamma/\delta$ (linha vertical)
indica o número de presas necessário para manter os predadores em equilíbrio. Quando as
presas estão abaixo desse limiar, os predadores declinam; acima, crescem.

\begin{lstlisting}[language=Python, caption=Simulacao e retrato de fase do Lotka-Volterra]
import numpy as np
from scipy.integrate import odeint
import matplotlib.pyplot as plt

def lotka_volterra(X, t, alpha, beta, delta, gamma):
    x, y = X
    dxdt = alpha*x - beta*x*y
    dydt = delta*x*y - gamma*y
    return [dxdt, dydt]

alpha, beta, delta, gamma = 1.0, 0.5, 0.5, 1.0
t = np.linspace(0, 50, 2000)

# Multiplas condicoes iniciais
fig, axes = plt.subplots(1, 2, figsize=(12, 5))

for x0, y0 in [(2, 1), (3, 1.5), (1.5, 0.5)]:
    sol = odeint(lotka_volterra, [x0, y0], t, args=(alpha, beta, delta, gamma))
    axes[0].plot(t, sol[:, 0], label=f'x(0)={x0}')
    axes[0].plot(t, sol[:, 1], '--')
    axes[1].plot(sol[:, 0], sol[:, 1])

# Equilibrio
axes[1].plot(gamma/delta, alpha/beta, 'ro', markersize=10, label='Equilibrio')
axes[0].set_xlabel('Tempo'); axes[0].set_ylabel('Populacao')
axes[1].set_xlabel('Presas (x)'); axes[1].set_ylabel('Predadores (y)')
plt.tight_layout(); plt.show()
\end{lstlisting}

\section{Análise de Estabilidade}

\subsection{Linearização e o Jacobiano}

A estabilidade de um ponto de equilíbrio $\mathbf{x}^*$ é analisada por linearização. Seja
$\boldsymbol{\xi} = \mathbf{x} - \mathbf{x}^*$ a perturbação em torno do equilíbrio. Para
$\boldsymbol{\xi}$ pequeno, a dinâmica é aproximada pelo sistema linear
\begin{equation}
  \frac{d\boldsymbol{\xi}}{dt} \approx J(\mathbf{x}^*)\,\boldsymbol{\xi},
\end{equation}
onde $J$ é o \emph{jacobiano} de $\mathbf{f}$ avaliado no equilíbrio:
$J_{ij} = \partial f_i / \partial x_j\big|_{\mathbf{x}^*}$.

A estabilidade é determinada pelos autovalores $\lambda_k$ de $J$: o equilíbrio é
\textbf{estável} se $\operatorname{Re}(\lambda_k) < 0$ para todo $k$, \textbf{instável}
se ao menos um autovalor tem parte real positiva, e \textbf{marginalmente estável} (análise
não-linear necessária) quando algum autovalor tem parte real nula.

Para sistemas bidimensionais, a equação característica de $J$ se reduz a
\begin{equation}
  \lambda^2 - \operatorname{tr}(J)\,\lambda + \det(J) = 0,
\end{equation}
com soluções $\lambda = \frac{1}{2}\left[\operatorname{tr}(J) \pm \sqrt{\operatorname{tr}(J)^2 - 4\det(J)}\right]$.

\subsection{Classificação dos Pontos de Equilíbrio}

A natureza do ponto fixo em 2D é completamente determinada por dois escalares:
$\tau = \operatorname{tr}(J)$ e $\Delta = \det(J)$. O diagrama $(\tau, \Delta)$
organiza todos os tipos possíveis:

\begin{itemize}
\item Se $\Delta < 0$: os autovalores são reais com sinais opostos. O ponto é um
  \textbf{ponto de sela} --- instável (um autovetor estável, um instável).

\item Se $\Delta > 0$ e $\tau^2 > 4\Delta$: autovalores reais de mesmo sinal.
  Se $\tau < 0$: \textbf{nó estável} (convergência exponencial, sem oscilações).
  Se $\tau > 0$: \textbf{nó instável} (divergência exponencial).

\item Se $\Delta > 0$ e $\tau^2 < 4\Delta$: autovalores complexos conjugados
  $\lambda = \tau/2 \pm i\omega$, com $\omega = \sqrt{4\Delta - \tau^2}/2$.
  Se $\tau < 0$: \textbf{espiral estável} (convergência oscilatória amortecida).
  Se $\tau > 0$: \textbf{espiral instável} (divergência oscilatória).
  Se $\tau = 0$: \textbf{centro} (órbitas fechadas conservativas).
\end{itemize}

Esse diagrama é uma ferramenta poderosa: basta calcular traço e determinante do jacobiano
para classificar o equilíbrio sem resolver explicitamente os autovalores.

\begin{figure}[htbp]
\centering
\begin{tikzpicture}[scale=1.0]
    % Diagrama de classificação
    \draw[->, thick] (-3.5,0) -- (3.5,0) node[right, font=\small] {$\tau$};
    \draw[->, thick] (0,-3) -- (0,3) node[above, font=\small] {$\Delta$};

    % Parábola discriminante tau² = 4 Delta
    \draw[thick, azulescuro, dashed, domain=-2.8:2.8, samples=100]
        plot (\x, {(\x*\x)/4});
    \node[right, font=\footnotesize, azulescuro] at (1.4, 0.5) {$\tau^2 = 4\Delta$};

    % Regiões coloridas
    \fill[red!10] (-3.5,0) rectangle (3.5,-3); % sela
    \fill[verdeacido!15] (-3.5,0) -- (-3.5,0.2) .. controls (-1.5,0.5) and (0,0.2) .. (3.5,0.3) -- (3.5,3) -- (-3.5,3) -- cycle; % acima da parábola
    \fill[azulclaro!25] (0,0) ellipse (1.5 and 0.4); % foco

    % Labels
    \node[red, font=\bfseries\large] at (-2, -1.5) {Sela};
    \node[verdeescuro, font=\bfseries] at (0, 2.3) {Espiral/Centro};
    \node[azulescuro, font=\bfseries] at (-2, 0.7) {Nó estável};
    \node[azulescuro, font=\bfseries] at (2, 0.7) {Nó instável};
    \node at (0, 0) {+};

\end{tikzpicture}
\caption{Diagrama de classificação dos pontos de equilíbrio em sistemas 2D pelo par $(\tau = \operatorname{tr}(J), \Delta = \det(J))$. A parábola $\tau^2 = 4\Delta$ separa regiões de autovalores reais (nós e selas) e complexos (espirais e centros). O eixo $\Delta < 0$ corresponde a pontos de sela. Para $\Delta > 0$, o sinal de $\tau$ determina se o equilíbrio é estável (esquerda) ou instável (direita).}
\label{fig:04-classificacao}
\end{figure}

\begin{lstlisting}[language=Python, caption=Calculo do Jacobiano e autovalores com SymPy]
import sympy as sp
import numpy as np

x, y = sp.symbols('x y')

# Sistema de exemplo: Lotka-Volterra com capacidade de suporte
alpha, beta, delta, gamma = 1.0, 0.5, 0.5, 1.0
f1 = x*(1 - x/10 - 0.5*y)
f2 = y*(-0.75 + 0.5*x)

# Jacobiano simbolico
J = sp.Matrix([[sp.diff(f1, x), sp.diff(f1, y)],
               [sp.diff(f2, x), sp.diff(f2, y)]])

# Encontrar equilibrios
equilibrios = sp.solve([f1, f2], [x, y])
print("Pontos de equilibrio:", equilibrios)

# Analisar estabilidade em cada equilibrio
for eq in equilibrios:
    x_star, y_star = float(eq[0]), float(eq[1])
    J_num = np.array(J.subs([(x, x_star), (y, y_star)])).astype(float)
    autovalores = np.linalg.eigvals(J_num)
    tau = np.trace(J_num)
    det = np.linalg.det(J_num)
    estavel = all(np.real(autovalores) < 0)
    print(f"\nEq ({x_star:.2f}, {y_star:.2f}): tr={tau:.3f}, det={det:.3f}")
    print(f"  Autovalores: {autovalores}")
    print(f"  Estavel: {estavel}")
\end{lstlisting}


\section{Bifurcações}

\begin{definicaoenv}[Bifurcação]
Uma bifurcação é uma mudança qualitativa na dinâmica de um sistema quando um parâmetro de
controle $\mu$ atravessa um valor crítico $\mu_c$. Novos equilíbrios aparecem ou
desaparecem, equilíbrios existentes trocam de estabilidade, ou surgem novas estruturas
dinâmicas como ciclos limite.
\end{definicaoenv}

As bifurcações são biologicamente fundamentais porque explicam transições abruptas ou
graduais entre estados funcionais distintos: a diferenciação celular (transição entre dois
fenótipos estáveis), o surgimento de ritmos circadianos (transição para oscilações), a
bistabilidade de switches genéticos (dois estados atratores separados por uma separatriz),
e o limiar de patologias (onde um estado saudável perde estabilidade).

\subsection{Bifurcação Saddle-Node}

A bifurcação saddle-node (ou \emph{fold}) é a forma mais comum de equilíbrios aparecerem
ou desaparecerem. Sua forma normal é
\begin{equation}
  \frac{dx}{dt} = \mu + x^2.
\end{equation}
Para $\mu < 0$, o sistema tem dois equilíbrios: $x^* = -\sqrt{-\mu}$ (estável) e
$x^* = +\sqrt{-\mu}$ (instável). Em $\mu = 0$, os dois equilíbrios colidem num único ponto
semiestável. Para $\mu > 0$, nenhum equilíbrio existe e todas as trajetórias divergem.

Na biologia, a bifurcação saddle-node aparece em modelos de \emph{bistabilidade}: circuitos
com feedback positivo que podem existir em dois estados estáveis (``ligado'' e ``desligado'')
separados por um limiar instável (ponto de sela). Uma perturbação suficientemente grande
pode empurrar o sistema de um estado para o outro de forma irreversível --- o que se conhece
como \emph{histerese}.

\subsection{Bifurcação Transcrítica}

A bifurcação transcrítica é característica de sistemas onde um estado trivial (como extinção
populacional, $x^* = 0$) sempre existe. Sua forma normal é
\begin{equation}
  \frac{dx}{dt} = \mu x - x^2.
\end{equation}
Para $\mu < 0$, $x^* = 0$ é estável e $x^* = \mu$ é instável. Em $\mu = 0$, os dois
ramos colidem e trocam de estabilidade: para $\mu > 0$, $x^* = 0$ torna-se instável e
$x^* = \mu > 0$ é o novo equilíbrio estável. O modelo logístico de crescimento populacional
exibe exatamente este comportamento: para $r < 1$ (recursos insuficientes), a extinção é o
único equilíbrio estável; para $r > 1$, a população sobrevive num equilíbrio positivo.

\subsection{Bifurcação Pitchfork}

A bifurcação pitchfork ocorre em sistemas com simetria e produz splitting de equilíbrios.
A forma supercrítica é
\begin{equation}
  \frac{dx}{dt} = \mu x - x^3.
\end{equation}
Para $\mu < 0$, apenas $x^* = 0$ existe (estável). Para $\mu > 0$, $x^* = 0$ torna-se
instável e surgem dois novos equilíbrios estáveis $x^* = \pm\sqrt{\mu}$ --- uma transição
suave e contínua. A forma subcrítica ($\dot{x} = \mu x + x^3$) produz comportamento
oposto: para $\mu < 0$, dois equilíbrios instáveis flanqueiam $x^* = 0$ estável; quando
$\mu > 0$, todos os equilíbrios desaparecem e o sistema diverge catastroficamente. A
pitchfork subcrítica é relevante em modelos de diferenciação celular, onde pequenas
perturbações podem induzir transições irreversíveis de estado.

\subsection{Bifurcação de Hopf}

A bifurcação de Hopf é o mecanismo gerador de oscilações biológicas. Um equilíbrio estável
perde estabilidade quando um par de autovalores complexos conjugados $\lambda = \alpha(\mu) \pm i\omega(\mu)$
atravessa o eixo imaginário: $\alpha(\mu_c) = 0$, com $d\alpha/d\mu|_{\mu_c} \neq 0$.

\begin{figure}[htbp]
\centering
\begin{tikzpicture}[scale=0.9]
    % Painel 1: antes da Hopf (espiral estável)
    \begin{scope}[xshift=0cm]
        \draw[->, thick] (-2,0) -- (2,0) node[right, font=\small] {$x$};
        \draw[->, thick] (0,-2) -- (0,2) node[above, font=\small] {$y$};
        \fill[azulescuro] (0,0) circle (2.5pt);
        \node[below right, font=\footnotesize, azulescuro] at (0,0) {$\mu < \mu_c$};
        % Espiral convergente
        \draw[thick, verdeacido, smooth, domain=0:540, samples=100]
            plot ({1.5*cos(\x)*exp(-\x/220)}, {1.5*sin(\x)*exp(-\x/220)});
        \node[below, font=\small] at (0,-2.5) {Espiral estável};
    \end{scope}

    % Painel 2: depois da Hopf (ciclo limite)
    \begin{scope}[xshift=5cm]
        \draw[->, thick] (-2,0) -- (2,0) node[right, font=\small] {$x$};
        \draw[->, thick] (0,-2) -- (0,2) node[above, font=\small] {$y$};
        \fill[red] (0,0) circle (2.5pt);
        \node[below right, font=\footnotesize, red] at (0,0) {$\mu > \mu_c$};
        % Ciclo limite
        \draw[thick, laranjacel, line width=1.5pt] (0,0) circle (1.2);
        % Seta no ciclo
        \draw[->, very thick, laranjacel] (1.0,0.5) -- (1.2,0);
        % Espiral divergente
        \draw[thick, roxodna, smooth, dashed, domain=0:540, samples=80]
            plot ({0.3*cos(\x)*exp(\x/600)}, {0.3*sin(\x)*exp(\x/600)});
        \node[below, font=\small] at (0,-2.5) {Ciclo limite estável};
    \end{scope}
\end{tikzpicture}
\caption{Bifurcação de Hopf supercrítica. (Esquerda) Para $\mu < \mu_c$, o equilíbrio é um foco estável: trajetórias vizinhas convergem para a origem em espiral amortecida. (Direita) Para $\mu > \mu_c$, o equilíbrio se torna instável e surge um \emph{ciclo limite} estável --- uma órbita fechada de raio $A \sim \sqrt{\mu - \mu_c}$ para a qual todas as trajetórias vizinhas convergem. Esta é a transição que gera oscilações auto-sustentadas em sistemas biológicos como ritmos circadianos e oscilações de NF-$\kappa$B.}
\label{fig:04-hopf}
\end{figure}

Na bifurcação de Hopf \emph{supercrítica}, para $\mu > \mu_c$, surge um ciclo limite
estável ao redor do equilíbrio (agora instável). A amplitude do ciclo cresce como
$A \sim \sqrt{\mu - \mu_c}$ --- uma transição suave. Na bifurcação \emph{subcrítica}, um
ciclo limite instável colapsa sobre o equilíbrio, que então perde estabilidade abruptamente.

O oscilador circadiano, as oscilações de NF-$\kappa$B em resposta a LPS, e as oscilações
glicolíticas são exemplos de sistemas biológicos que passam por bifurcações de Hopf quando
parâmetros como taxas de síntese ou degradação são variados.

\begin{lstlisting}[language=Python, caption=Diagrama de bifurcacao via varredura de parametro]
import numpy as np
import matplotlib.pyplot as plt
from scipy.integrate import odeint

def hopf_normal(X, t, mu):
    x, y = X
    r2 = x**2 + y**2
    dxdt = mu*x - y - x*r2
    dydt = x + mu*y - y*r2
    return [dxdt, dydt]

t = np.linspace(0, 100, 5000)
mu_values = np.linspace(-0.5, 1.0, 30)
amplitudes = []

for mu in mu_values:
    sol = odeint(hopf_normal, [0.1, 0.1], t, args=(mu,))
    # Amplitude no regime estacionario (ultimos 20%)
    x_ss = sol[int(0.8*len(t)):, 0]
    amplitudes.append(np.max(x_ss) - np.min(x_ss))

plt.figure(figsize=(8, 5))
plt.plot(mu_values, amplitudes, 'b-o', markersize=4)
plt.axvline(0, color='r', linestyle='--', label='Bifurcacao de Hopf')
plt.xlabel('Parametro mu'); plt.ylabel('Amplitude do ciclo')
plt.title('Bifurcacao de Hopf supercritica')
plt.legend(); plt.grid(True); plt.show()
\end{lstlisting}

\section{Oscilações e Ciclos Limite}

\begin{definicaoenv}[Ciclo Limite]
Um ciclo limite é uma órbita periódica \emph{isolada} no espaço de fase, para a qual
trajetórias próximas convergem (ciclo limite \emph{estável}) ou divergem (ciclo limite
\emph{instável}). Ao contrário do \emph{centro} do sistema linear conservativo --- onde
órbitas formam uma família infinita dependente das condições iniciais --- o ciclo limite tem
amplitude e período fixos, determinados pela não-linearidade do sistema.
\end{definicaoenv}

Ciclos limite são essenciais para modelar ritmos biológicos robustos: o relógio circadiano
de uma célula continua a oscilar com período de $\sim$24 horas mesmo após perturbações de
temperatura ou luz, porque o ciclo limite tem amplitude e período intrínsecos, independentes
das condições iniciais.

\subsection{O Teorema de Poincaré-Bendixson}

Em sistemas bidimensionais, a existência de ciclos limite é garantida pelo teorema de
Poincaré-Bendixson: se existe uma região limitada $R$, positivamente invariante (trajetórias
não saem de $R$) e sem pontos de equilíbrio, então qualquer trajetória que permanece em $R$
converge para um ciclo limite periódico. Este teorema é uma ferramenta poderosa para provar
existência de oscilações sem encontrar a solução explícita. Ele também implica que
\emph{caos não é possível em sistemas autônomos bidimensionais} --- para haver caos
determinístico, são necessárias pelo menos três dimensões.

\subsection{O Oscilador de Van der Pol}

O oscilador de Van der Pol é o paradigma dos osciladores não-lineares com ciclo limite. Na
forma de sistema de primeira ordem:
\begin{align}
  \frac{dx}{dt} &= y \\
  \frac{dy}{dt} &= \mu(1 - x^2)y - x,
\end{align}
onde $\mu > 0$ controla a não-linearidade. Para $\mu$ pequeno, o oscilador é quase
harmônico, com amplitude e frequência próximas de 1. Para $\mu$ grande, o sistema exibe
\emph{oscilações de relaxação}: passa longos períodos em estados quase estáticos,
interrompidos por transições rápidas --- o comportamento típico de osciladores biológicos
com dinâmica rápido-lento, como o potencial de ação cardíaco.

A análise de estabilidade da origem revela que o jacobiano avaliado em $(0,0)$ tem traço
$\operatorname{tr}(J) = \mu > 0$ e $\det(J) = 1 > 0$, portanto a origem é sempre uma
espiral instável para $\mu > 0$. A existência de um ciclo limite estável --- provável pelo
teorema de Poincaré-Bendixson, aplicado a uma região anular adequada --- é o que garante
que o sistema não diverge: as trajetórias se estabilizam no ciclo.

\subsection{Osciladores Biológicos}

Os sistemas biológicos possuem osciladores em todas as escalas de tempo. Os
\emph{ritmos circadianos} operam com período de $\sim$24 horas e são gerados por loops de
feedback transcricional negativo: a proteína PER inibe seu próprio gene, criando um atraso
temporal que produz oscilações autossustentadas. As \emph{oscilações de cálcio} intracelular
têm período de segundos a minutos e mediam sinalização em neurônios, células imunes e
oócitos. As \emph{oscilações glicolíticas} em leveduras (período $\sim$5 min) são geradas
pela ativação alostérica da fosfofrutoquinase (PFK) pelo ADP --- um loop de feedback
positivo balanceado por consumo de ATP. O \emph{ciclo celular} coordena duplicação de DNA e
divisão celular com período de horas, combinando ciclos limite (para a progressão periódica)
com switches biestáveis (para as transições irreversíveis de ponto de controle).

\begin{lstlisting}[language=Python, caption=Oscilador de Van der Pol e Goodwin]
import numpy as np
from scipy.integrate import odeint
import matplotlib.pyplot as plt

# Van der Pol
def van_der_pol(X, t, mu):
    x, y = X
    return [y, mu*(1 - x**2)*y - x]

# Oscilador de Goodwin (modelo de relogio biologico)
def goodwin(X, t, n=4, k1=1, k2=1, k3=1):
    x, y, z = X
    return [
        1/(1 + z**n) - k1*x,
        x - k2*y,
        y - k3*z
    ]

t = np.linspace(0, 50, 2000)
fig, axes = plt.subplots(2, 2, figsize=(12, 8))

# Van der Pol para diferentes mu
for i, mu in enumerate([0.5, 2.0]):
    sol = odeint(van_der_pol, [0.1, 0.1], t, args=(mu,))
    axes[0, i].plot(t, sol[:, 0])
    axes[0, i].set_title(f'Van der Pol, mu={mu}')
    axes[0, i].set_xlabel('Tempo'); axes[0, i].set_ylabel('x(t)')

# Goodwin
sol_gw = odeint(goodwin, [0.5, 0.5, 0.5], t)
axes[1, 0].plot(t, sol_gw)
axes[1, 0].set_title('Oscilador de Goodwin')
axes[1, 0].legend(['x', 'y', 'z'])
axes[1, 1].plot(sol_gw[:, 0], sol_gw[:, 2])
axes[1, 1].set_title('Retrato de fase Goodwin (x vs z)')

plt.tight_layout(); plt.show()
\end{lstlisting}


\section{Análise de Sensibilidade}

A análise de sensibilidade quantifica como as saídas do modelo respondem a variações nos
parâmetros --- pergunta fundamental tanto para validação de modelos quanto para identificar
alvos terapêuticos. A distinção central é entre análise \emph{local} (derivadas parciais
num ponto nominal) e \emph{global} (exploração de todo o espaço paramétrico).

\subsection{Sensibilidade Local}

O \emph{coeficiente de sensibilidade local} de uma variável $x_i$ em relação a um
parâmetro $\theta_j$ é a derivada parcial normalizada:
\begin{equation}
  S_{ij} = \frac{\partial \ln x_i}{\partial \ln \theta_j}
           = \frac{\theta_j}{x_i}\,\frac{\partial x_i}{\partial \theta_j}.
\end{equation}
Um valor $S_{ij} = 2$ significa que um aumento de 1\% em $\theta_j$ produz um aumento de
2\% em $x_i$. Coeficientes próximos de zero indicam parâmetros aos quais o modelo é
robusto; valores grandes indicam vulnerabilidades ou, na perspectiva terapêutica, alvos de
intervenção promissores.

Para sistemas de EDOs, a sensibilidade é uma função do tempo $S_{ij}(t)$, que satisfaz a
\emph{equação de sensibilidade} obtida diferenciando o sistema em relação a $\theta_j$:
\begin{equation}
  \frac{d}{dt}\!\left(\frac{\partial x_i}{\partial \theta_j}\right)
  = \sum_k \frac{\partial f_i}{\partial x_k}\,\frac{\partial x_k}{\partial \theta_j}
    + \frac{\partial f_i}{\partial \theta_j}.
\end{equation}
Em forma matricial, as equações de sensibilidade formam um sistema linear forçado que deve
ser integrado \emph{juntamente} com o sistema original.

\subsection{Sensibilidade Global: Índices de Sobol}

Para incerteza paramétrica ampla, a \emph{análise de sensibilidade global} explora todo o
espaço de parâmetros e decompõe a variância total da saída $Y$ em contribuições de cada
parâmetro. Os \emph{índices de Sobol} quantificam isso:
\begin{equation}
  S_i = \frac{V[\mathbb{E}(Y|p_i)]}{V[Y]}, \quad
  S_i^T = \frac{\mathbb{E}[V(Y|p_{\sim i})]}{V[Y]},
\end{equation}
onde $S_i$ é o índice de primeira ordem (efeito principal de $p_i$) e $S_i^T$ é o índice
total (incluindo interações com outros parâmetros). A diferença $S_i^T - S_i$ mede a
importância das interações de $p_i$ com outros parâmetros.

A amostragem do espaço paramétrico emprega técnicas como \emph{Latin Hypercube Sampling}
(LHS) --- que estratifica cada dimensão em $n$ intervalos iguais e amostra um ponto por
intervalo, garantindo cobertura uniforme --- ou \emph{sequências de Sobol}, que são
quasi-aleatórias e têm menor discrepância que amostragem de Monte Carlo puro.

\begin{lstlisting}[language=Python, caption=Analise de sensibilidade local e global]
import numpy as np
from scipy.integrate import odeint

# Sensibilidade local por perturbacao finita
def calc_sens(param_name, delta=0.01):
    def model(x, t, Vmax, Km, kdeg, S=8):
        return (Vmax*S)/(Km+S) - kdeg*x
    t = np.linspace(0, 50, 500)
    p0 = {'Vmax': 10, 'Km': 5, 'kdeg': 0.5}
    # Valor nominal
    y0 = odeint(model, 0, t, args=(p0['Vmax'], p0['Km'], p0['kdeg']))[-1, 0]
    # Valor perturbado
    p1 = p0.copy(); p1[param_name] *= (1 + delta)
    y1 = odeint(model, 0, t, args=(p1['Vmax'], p1['Km'], p1['kdeg']))[-1, 0]
    return ((y1 - y0)/y0) / delta  # sensibilidade normalizada

for p in ['Vmax', 'Km', 'kdeg']:
    s = calc_sens(p)
    print(f"S(x*, {p}) = {s:.3f}")

# Analise global com SALib (instalar: pip install SALib)
# from SALib.sample import saltelli
# from SALib.analyze import sobol
# problem = {'num_vars': 3, 'names': ['Vmax', 'Km', 'kdeg'],
#            'bounds': [[5, 15], [2.5, 7.5], [0.25, 0.75]]}
# param_values = saltelli.sample(problem, 1024)
# Y = np.array([(V*8)/(k*(K+8)) for V, K, k in param_values])
# Si = sobol.analyze(problem, Y)
# print("Indices S1:", dict(zip(problem['names'], Si['S1'])))
\end{lstlisting}


\section{Exercícios}

\begin{exercicio}
Simule o mapa logístico $x_{t+1} = r\,x_t(1 - x_t)$ para $r \in \{2.5, 3.2, 3.5, 3.9\}$
e $x_0 = 0.3$. (a) Plote as 100 primeiras iterações para cada valor de $r$. (b) Construa
o diagrama de bifurcação para $r \in [2.5, 4.0]$. (c) Para $r = 3.9$, compare trajetórias
iniciadas em $x_0 = 0.3$ e $x_0 = 0.30001$ --- quantas iterações são necessárias para
que divirjam visivelmente? O que esse experimento demonstra sobre o caos determinístico?
\end{exercicio}

\begin{exercicio}
Considere um canal iônico com taxa de abertura $\alpha = 0.3$ e taxa de fechamento
$\beta = 0.7$ por passo de tempo. (a) Escreva a matriz de transição $\mathbf{P}$.
(b) Calcule analiticamente a distribuição estacionária $\pi_A$ e $\pi_F$.
(c) Simule a cadeia por 1000 passos a partir do estado Fechado e verifique que a
fração de tempo em cada estado converge para a distribuição estacionária.
\end{exercicio}

\begin{exercicio}
Para o sistema de Lotka-Volterra com $\alpha = 0.4$, $\beta = 0.02$, $\delta = 0.01$,
$\gamma = 0.2$: (a) calcule o ponto de equilíbrio não-trivial $(x^*, y^*)$; (b) calcule
o jacobiano nesse ponto e determine os autovalores; (c) classifique o equilíbrio no
diagrama $(\tau, \Delta)$; (d) simule o sistema com \texttt{scipy.integrate.odeint} e
plote o retrato de fase para três condições iniciais distintas.
\end{exercicio}

\begin{exercicio}
Analise a estabilidade de todos os pontos de equilíbrio do sistema:
\begin{align*}
  \frac{dx}{dt} &= y \\
  \frac{dy}{dt} &= -x - y + x^3
\end{align*}
Calcule o jacobiano em cada equilíbrio, determine os autovalores e classifique o tipo
(nó, espiral, sela, centro). Esboce o retrato de fase qualitativamente.
\end{exercicio}

\begin{exercicio}
Para a equação normal da bifurcação pitchfork supercrítica $\dot{x} = \mu x - x^3$:
(a) encontre analiticamente todos os equilíbrios para $\mu < 0$, $\mu = 0$ e $\mu > 0$;
(b) determine a estabilidade de cada equilíbrio via derivada de $f$; (c) construa
computacionalmente o diagrama de bifurcação completo, desenhando ramos estáveis (linha
sólida) e instáveis (tracejado).
\end{exercicio}

\begin{exercicio}
Verifique numericamente a bifurcação de Hopf supercrítica no sistema:
\begin{align*}
  \frac{dx}{dt} &= \mu x - y - x(x^2 + y^2) \\
  \frac{dy}{dt} &= x + \mu y - y(x^2 + y^2)
\end{align*}
Para cada valor $\mu \in \{-0.5, -0.1, 0.1, 0.5, 1.0\}$: simule o sistema por tempo
suficiente, plote o retrato de fase e mida a amplitude do regime estacionário. Construa
um gráfico de amplitude vs.\ $\mu$ e compare com a previsão teórica $A \sim \sqrt{\mu}$.
\end{exercicio}

\begin{exercicio}
Implemente o oscilador de Van der Pol para $\mu \in \{0.1, 1, 5\}$ e: (a) compare as
formas de onda $x(t)$ --- qual delas é mais próxima de um seno? qual apresenta
oscilações de relaxação? (b) estime o período de oscilação em cada caso; (c) plote os
retratos de fase e identifique o ciclo limite.
\end{exercicio}

\begin{exercicio}
Implemente o oscilador de Goodwin com $n = 4$, $k_1 = k_2 = k_3 = 1$:
\begin{align*}
  \dot{x} &= \frac{1}{1+z^n} - k_1 x \\
  \dot{y} &= x - k_2 y \\
  \dot{z} &= y - k_3 z
\end{align*}
(a) Simule e verifique que o sistema exibe oscilações sustentadas. (b) Calcule os
coeficientes de sensibilidade local $S_{A, n}$, $S_{A, k_1}$, $S_{A, k_2}$ da amplitude
$A$ em relação a cada parâmetro. (c) Qual parâmetro controla mais fortemente a amplitude?
E o período?
\end{exercicio}

\begin{exercicio}[Projeto Integrador]
Escolha um modelo de rede regulatória bidimensional: o \emph{toggle switch} de Gardner
et al.\ (2000) ou o \emph{repressilator} de Elowitz e Leibler (2000).
(a) Implemente o modelo em Python. (b) Realize análise de estabilidade completa --- encontre
todos os equilíbrios e classifique-os. (c) Identifique uma bifurcação variando um parâmetro
chave (ex: força do feedback, taxa de degradação). (d) Conduza análise de sensibilidade
global (índices de Sobol) para determinar quais parâmetros controlam mais fortemente o
comportamento do sistema. (e) Interprete biologicamente: o que a topologia matemática
revela sobre a função do circuito?
\end{exercicio}

\section{Notas Bibliográficas}

O artigo seminal de May (1976) \cite{may1976}, \textit{Simple mathematical models with very
complicated dynamics}, publicado na \textit{Nature}, é leitura obrigatória e acessível ---
introduziu o conceito de caos determinístico a uma audiência científica ampla. A constante
de Feigenbaum e a universalidade da duplicação de período são detalhadas em Feigenbaum
(1978) \cite{feigenbaum1978}.

A teoria das bifurcações é tratada com rigor e clareza em Strogatz (2015)
\cite{strogatz2015}, \textit{Nonlinear Dynamics and Chaos} --- o texto introdutório mais
recomendado para biólogos com formação matemática básica. O livro de Murray (2002)
\cite{murray2002}, \textit{Mathematical Biology}, cobre desde crescimento populacional até
eletrofisiologia cardíaca. Keener e Sneyd (2009) \cite{keener2009}, \textit{Mathematical
Physiology}, são a referência para aplicações biomédicas detalhadas.

Para osciladores biológicos, Novák e Tyson (2008) \cite{novak2008} fornecem uma revisão
exemplar dos princípios de design de osciladores bioquímicos. Tyson et al.\ (2003)
\cite{tyson2003} analisam o ciclo celular como sistema dinâmico. O artigo de Ferrell (2012)
\cite{ferrell2012} conecta bistabilidade e bifurcações ao conceito de ``paisagem epigenética
de Waddington'', fundamental para biologia do desenvolvimento.

A análise de sensibilidade global é detalhada em Saltelli et al.\ (2008)
\cite{saltelli2008}, \textit{Global Sensitivity Analysis: The Primer}. A biblioteca SALib
para Python implementa todos os métodos discutidos (github.com/SALib/SALib).

\begin{thebibliography}{99}

\bibitem{may1976}
May R.M. (1976).
Simple mathematical models with very complicated dynamics.
\textit{Nature} \textbf{261}: 459--467.

\bibitem{feigenbaum1978}
Feigenbaum M.J. (1978).
Quantitative universality for a class of nonlinear transformations.
\textit{Journal of Statistical Physics} \textbf{19}: 25--52.

\bibitem{strogatz2015}
Strogatz S.H. (2015).
\textit{Nonlinear Dynamics and Chaos: With Applications to Physics, Biology, Chemistry,
and Engineering}, 2\textordfeminine\ ed.
Westview Press.

\bibitem{murray2002}
Murray J.D. (2002).
\textit{Mathematical Biology I: An Introduction}, 3\textordfeminine\ ed.
Springer.

\bibitem{keener2009}
Keener J., Sneyd J. (2009).
\textit{Mathematical Physiology}, 2\textordfeminine\ ed.
Springer.

\bibitem{novak2008}
Novák B., Tyson J.J. (2008).
Design principles of biochemical oscillators.
\textit{Nature Reviews Molecular Cell Biology} \textbf{9}: 981--991.

\bibitem{tyson2003}
Tyson J.J., Chen K.C., Novák B. (2003).
Sniffers, buzzers, toggles and blinkers: dynamics of regulatory and signaling pathways
in the cell.
\textit{Current Opinion in Cell Biology} \textbf{15}: 221--231.

\bibitem{ferrell2012}
Ferrell J.E. (2012).
Bistability, bifurcations, and Waddington's epigenetic landscape.
\textit{Current Biology} \textbf{22}: R458--R466.

\bibitem{saltelli2008}
Saltelli A. et al. (2008).
\textit{Global Sensitivity Analysis: The Primer}.
Wiley.

\end{thebibliography}
